\section{Theoretical Calculations}

\subsection{Notation}

在此处统一定义符号、变量与假设。例如，令
\begin{equation}
    \bm{x} = (x_1, x_2, \dots, x_n)^\top \in \mathbb{R}^n .
\end{equation}

\subsection{Model Assumptions}

在此处列出理论模型的基本假设。

\begin{definition}
给定目标函数 $f:\mathbb{R}^n \to \mathbb{R}$，若 $f$ 二阶连续可微，则其一阶与二阶局部展开可写为
\begin{equation}
    f(\bm{x} + \Delta \bm{x})
    =
    f(\bm{x})
    + \nabla f(\bm{x})^\top \Delta \bm{x}
    + \frac{1}{2}\Delta \bm{x}^\top \nabla^2 f(\bm{x}) \Delta \bm{x}
    + o(\|\Delta \bm{x}\|^2).
\end{equation}
\end{definition}

\subsection{Derivation}

在此处逐步写入核心推导。

\begin{align}
    A &= B + C, \\
    \frac{\partial A}{\partial x_i}
      &= \frac{\partial B}{\partial x_i}
       + \frac{\partial C}{\partial x_i}.
\end{align}

\subsection{Conclusion}

在此处总结理论计算得到的关键结论，并标记后续需要证明或验证的部分。
